% A comparison sheet followed by all three full-size concept proofs.
\documentclass[11pt,a4paper,landscape]{article}
\usepackage[a4paper,landscape,margin=16mm]{geometry}
% Shared font contract for the interior and standalone cover.
% Resolve via kpathsea and pass explicit directories: otherwise luaotfload may
% silently prefer a different same-named font in the user's Library/Fonts.
\usepackage{fontspec}
\defaultfontfeatures{Ligatures=TeX}
\edef\irserifpath{\directlua{
  local p = kpse.find_file("LibertinusSerif-Regular.otf", "opentype fonts")
  if p then tex.sprint((p:gsub("[^/]+$", "")))
  else tex.error("TeX-distributed Libertinus Serif fonts were not found") end
}}
\edef\irsanspath{\directlua{
  local p = kpse.find_file("SourceSansPro-Regular.otf", "opentype fonts")
  if p then tex.sprint((p:gsub("[^/]+$", "")))
  else tex.error("TeX-distributed Source Sans Pro fonts were not found") end
}}
\setmainfont{LibertinusSerif-Regular.otf}[
  Path=\irserifpath,
  ItalicFont=LibertinusSerif-Italic.otf,
  BoldFont=LibertinusSerif-Bold.otf,
  BoldItalicFont=LibertinusSerif-BoldItalic.otf,
  Numbers={Lining,Proportional}
]
\setsansfont{SourceSansPro-Regular.otf}[
  Path=\irsanspath,
  ItalicFont=SourceSansPro-RegularIt.otf,
  BoldFont=SourceSansPro-Semibold.otf,
  BoldItalicFont=SourceSansPro-SemiboldIt.otf,
  Scale=MatchLowercase,
  Numbers={Lining,Proportional}
]
% Semantic IPA tag: inherit the surrounding body font, size, and shape.
\newcommand{\ipagloss}[1]{{#1}}

\usepackage{graphicx,xcolor,pdfpages}
\usepackage[hidelinks]{hyperref}
\definecolor{reviewink}{HTML}{102F3A}
\hypersetup{pdftitle={Iridian - three cover directions},pdfauthor={Roel Christian Yambao}}
\pagestyle{empty}
\setlength{\parindent}{0pt}
\begin{document}
{\color{reviewink}\sffamily\fontsize{23}{28}\selectfont Iridian / Three cover directions}\par
\vspace{2mm}
{\sffamily\small All three preserve the original sideways branching. Full-size proofs follow.}\par
\vspace{6mm}
\begin{minipage}[t]{0.31\textwidth}
  \centering\includegraphics[page=1,height=105mm]{build/cover/concepts.pdf}\par
  \vspace{4mm}\raggedright
  {\sffamily\bfseries 01\quad Constituent structure}\par
  \vspace{1.5mm}
  {\small\textbf{Recommended.} Deep teal, warm ivory and an open serif title. The branching silhouette presents grammar as an organised system while retaining the original cover's identity.}
\end{minipage}\hfill
\begin{minipage}[t]{0.31\textwidth}
  \centering\includegraphics[page=2,height=105mm]{build/cover/concepts.pdf}\par
  \vspace{4mm}\raggedright
  {\sffamily\bfseries 02\quad The analytical margin}\par
  \vspace{1.5mm}
  {\small An ivory field meets a dark outer margin; the tree changes contrast as it crosses the boundary. The composition evokes the relation between a text and its analytical apparatus.}
\end{minipage}\hfill
\begin{minipage}[t]{0.31\textwidth}
  \centering\includegraphics[page=3,height=105mm]{build/cover/concepts.pdf}\par
  \vspace{4mm}\raggedright
  {\sffamily\bfseries 03\quad Descriptive strata}\par
  \vspace{1.5mm}
  {\small Warm paper above an oxblood field separates the bibliographic identity from the structural artwork. Its horizontal division gives the reference volume a quieter, more classical character.}
\end{minipage}
\includepdf[pages=-,fitpaper=true]{build/cover/concepts.pdf}
\end{document}
